\BourbakiExerciseGroup

\begin{enumerate}
\def\labelenumi{\arabic{enumi})}
\tightlist
\item
称点 \(a \in \mathbb{R}\) 为子集 \(\mathbf{A}\) （ \(\mathbb{R}\) 的）的左附着点，如果它属于 \(\mathbf{A}\) 的闭包，并且存在一个区间 \(]a, b[ (a < b)\) ，它不含 \(\mathbf{A}\) 的任何点。试证：子集 \(\mathbf{A}\) （ \(\mathbb{R}\) 的）的左附着点所成的集是可数的（在左附着点所成的集与一个由两两不相交的开区间组成的集之间建立一一对应）。由此试证 \(\mathbb{R}\) 的每个良序子集都是可数的。
\end{enumerate}

¶ 2) 设 \(\mathbf{A}\) 为 \(\mathbb{R}\) 中一个可数稠密子集，则它不是闭集。{[}若 \((a_{n})\) 是把 \(\mathbf{A}\) 的点按某个次序排列所得的序列，试定义一个区间序列 \([b_{n}, c_{n}]\) ，使得 \(b_{n-1} < b_{n} < c_{n} < c_{n-1}\) 对一切 \(n\) 成立，并且使得 \([b_{n}, c_{n}]\) 不含任何点 \(a_{k}\) ，其指标满足 \(k \leqslant n\) ；然后利用定理 2.{]}由此推出 \(\mathbb{R}\) 是不可数的（参看 § 8，第 6 节）。

\begin{enumerate}
\def\labelenumi{\arabic{enumi})}
\setcounter{enumi}{2}
\item
设 \((\mathbf{I}_n)\) 是 \(\mathbf{R}\) 中非空开区间的一个无穷序列，满足 \(\overline{\mathbf{I}}_n \cap \overline{\mathbf{I}}_m = \emptyset\) （ \(m \neq n\) ）。试证：诸 \(\mathbf{I}_n\) 的并的余集是一个完备集（第一章，§ 2，第 6 节）；特别地，Cantor 三分集是完备集。
\item
与 Cantor 三分集邻接的各区间的长度之和等于 1。试定义 {[}0, 1{]} 的一个全不连通完备子集，使得与 A 邻接且含于 {[}0, 1{]} 内的各区间的长度之和等于给定的数 m，它满足 \(0 < m \leqslant 1\) 。
\item
设 \(l\) 是一个数 \(> 0\) ，并假设对每个点 \(x \in \mathbf{R}\) 对应着一个开区间 \(\mathbf{I}(x)\) ，它以 \(x\) 为中心，长度 \(\leqslant l\) 。试证：每个紧区间 \([a, b]\) 都可被有限多个区间 \(\mathbf{I}(x_i)\) 覆盖，这些区间的长度之和 \(\leqslant l + 2(b - a)\) 。（试证：若命题对每个区间 \([a, x]\) （满足 \(a \leqslant x < c\) ）为真，则存在 \(d > c\) ，使它对每个区间 \([a, y]\) （满足 \(a \leqslant y < d\) ）为真。）若 \(l = (b - a)/n\) ，其中 \(n\) 是一个整数 \(\geqslant 1\) ，试证该结果不能再改进。
\end{enumerate}

¶ 6) a) 设 X 是一个非空全序集，赋予拓扑 \(\mathcal{C}_{0}(X)\) （第一章，§ 2，习题 5）。试证：X 是紧的当且仅当 X 的每个子集都有上确界，换言之，当且仅当 X 是一个完备格（集合论，第三章，§ 1，习题 11）。（为证条件必要，可按定理 3 的方式论证；为证充分性，考虑 X 上的一个滤子 \(\mathfrak{F}\) ，并证明：若 A 是 \(\mathfrak{F}\) 中诸集的下确界所成的集，则 A 的上确界是 \(\mathfrak{F}\) 的一个聚点。）

\begin{enumerate}
\def\labelenumi{\alph{enumi})}
\setcounter{enumi}{1}
\tightlist
\item
试给出一个完备格 X 的例子，它在拓扑 \(\mathcal{C}_{0}(\mathbf{X})\) 下不是紧的。
\end{enumerate}

¶ 7) 设 X 是一个非空全序集，赋予拓扑 \(\mathfrak{G}_{\mathfrak{0}}(\mathbf{X})\) 。

\begin{enumerate}
\def\labelenumi{\alph{enumi})}
\tightlist
\item
若 X 是连通的，试证它具有下列两个性质：
\end{enumerate}

α) X 的每个非空且上有界的子集都有上确界（若 A 非空且上有界，B 是 A 的上界所成的集，C 是 B 的下界所成的集，试证 \(B \cup C = X\) ，并且 B 与 C 都是闭集）。

β) 集 X 没有间隙，换言之（集合论，第三章，§ 1，习题 18），X 的每个开区间 {]}a, b{[} （a \textless{} b）都非空（利用定理 4 的论证）。

\begin{enumerate}
\def\labelenumi{\alph{enumi})}
\setcounter{enumi}{1}
\item
反之，试证若 X 具有性质 \(\alpha\) 与 \(\beta\) ），则 X 是连通的。（先证明每个闭区间 \([a, b]\) 都是连通的：设该区间是两个非空闭集 A 与 B 的不交并，且 \(a \in A\) ，考察 B 的下确界，由此得出矛盾。）
\item
试证：若 X 连通，则它是局部紧且局部连通的，并且 X 的仅有的连通子集就是区间（有界或无界）。
\end{enumerate}

\begin{enumerate}
\def\labelenumi{\arabic{enumi})}
\setcounter{enumi}{7}
\tightlist
\item
设 X 与 Y 是两个全序集，分别赋予拓扑 \(\mathcal{G}_{0}(X)\) 与 \(\mathcal{G}_{0}(Y)\) 。
\end{enumerate}

\begin{enumerate}
\def\labelenumi{\alph{enumi})}
\tightlist
\item
设 \(\mathbf{X}\) 连通，且设 \(f: \mathbf{X} \to \mathbf{Y}\) 是一个连续映射。试证：设 \(x, y\) 为 \(\mathbf{X}\) 中满足 \(x < y\) 的任意两点，则每个 \(z \in \mathbf{Y}\) ，只要它属于以 \(f(x)\) 与 \(f(y)\) 为端点的闭区间，就属于 \(f(\mathbf{X})\)
\end{enumerate}

（利用习题 7）。由此试证：\(f\) 是 \(\mathbf{X}\) 到 \(f(\mathbf{X})\) 上的同胚，当且仅当 \(f\) 连续且严格单调。

\begin{enumerate}
\def\labelenumi{\alph{enumi})}
\setcounter{enumi}{1}
\tightlist
\item
试给出有理直线 Q 到其自身的一个同胚的例子，它不是严格单调的。
\end{enumerate}

\begin{enumerate}
\def\labelenumi{\arabic{enumi})}
\setcounter{enumi}{8}
\tightlist
\item
  \begin{enumerate}
  \def\labelenumii{\alph{enumii})}
  \tightlist
  \item
设 X 是一个无间隙的可数全序集（习题 7）。试证：存在一个严格递增的双射映射 \(\varphi\) ，它把 X 映到以 o 与 i 为端点的 Q 的四个区间之一上。{[}把 X 的元素与 Q 的适当区间的点排成序列 \((a_n), (b_n)\) ，并用归纳法定义 \(\varphi\) 。{]}若 X 赋予拓扑 \(\tilde{C}_0(X)\) ，则 \(\varphi\) 是 X 到 \(\varphi(X)\) 上的同胚。
  \end{enumerate}
\end{enumerate}

\begin{enumerate}
\def\labelenumi{\alph{enumi})}
\setcounter{enumi}{1}
\item
试由 a) 推出：R 的开区间 {]}a, b{[} 的每个可数稠密子集都同胚于 Q。
\item
试由 a) 推出：若 X 是任一赋予拓扑 \(\mathfrak{G}_{0}(\mathbf{X})\) 的可数全序集，则存在一个严格递增的同胚，把 X 映到 Q 的一个子空间上{[}把 X 嵌入一个无间隙的可数全序空间 \(X'\) ，使得 \(\mathfrak{G}_{0}(\mathbf{X})\) 由 \(\mathfrak{G}_{0}(\mathbf{X}')\) 诱导{]}。
\end{enumerate}

\begin{enumerate}
\def\labelenumi{\arabic{enumi})}
\setcounter{enumi}{9}
\tightlist
\item
试证：Cantor 三分集 K 的每个可数子集，只要它在 K 中稠密且不含与 K 邻接的区间的任何端点，就同胚于有理直线 Q（利用习题 9）。
\end{enumerate}

¶ 11) a) 设 X 是带拓扑 \(\mathfrak{G}_{0}(\mathbf{X})\) 的全序集。若 X 连通且有一个可数稠密子集 A，试证：存在一个（由习题 8 知为严格单调的）同胚，把 X 映到以 o, i 为端点的 R 的区间之一上，并且把 A 映到 Q 与该区间的交上（利用习题 9 与 7）。

\begin{enumerate}
\def\labelenumi{\alph{enumi})}
\setcounter{enumi}{1}
\item
试由 a) 推出：若 B 是 R 的一个子集，它的余集在 R 中稠密，则存在 R 到其自身的同胚，它把 B 映到无理数集 CQ 的一个子集上（考虑一个含于 CB 中的 R 的可数稠密子集 A）。
\item
在按字典序排成全序的集合 \(\mathbf{R} \times \{0, 1\}\) 中，用 \(\mathbf{R}'\) 表示 \(\mathbf{Q} \times \{1\}\) 的余集。试证：\(\mathbf{R}'\) 赋予拓扑 \(\mathfrak{G}_0(\mathbf{R}')\) 后是局部紧且全不连通的，并且含有一个可数稠密子集 \(\mathbf{Q}'\) ，它同胚于 \(\mathbf{Q}\) ，但在 \(\mathbf{R}'\) 的非空开区间上诱导的拓扑不具有可数基；特别地，这样的区间不能同胚于 \(\mathbf{R}\) 的任何子集。
\end{enumerate}

¶ 12) a) 设 A 是一个良序集，用 I 表示 R 的区间 {[}0, 1{[} 。试证：按字典序排成全序并赋予拓扑 \(\mathfrak{G}_{0}(X)\) 的集合 X = A × I 是连通的（习题 7）；因此存在在拓扑 \(\mathfrak{G}_{0}(\mathbf{X})\) 下连通且具有任意势的全序集 X{[}参看 § 4，习题 7 b){]}。

\begin{enumerate}
\def\labelenumi{\alph{enumi})}
\setcounter{enumi}{1}
\tightlist
\item
取 A 为一个不可数的良序集，其中每个线段 \(\leftarrow, t]\) 都是可数的：相应的拓扑空间 X 称为 Alexandroff 半直线。试证：对 X 的每个区间 \([a, b]\) （a \textless{} b），存在一个严格递增的同胚，把 \([a, b]\) 映到 R 的区间 \([o, i]\) 上。（用超限归纳法证明：存在一个严格递增的同胚，把 \(A \cap [a, b]\) （把 A 等同于 X 中点 \((t, o)\) 所成之集）映到 \([o, i]\) 的一个子空间上。）
\end{enumerate}

¶ 13) 设 \(a\) 是一个无理数 \(>0\) 。对每个有理数 \(x\) ，令 \(f_{a}(x)\) 为区间 \([0, a]\) 中使 \(x - f_{a}(x)\) 是 \(a\) 的整数倍的那个实数。试证：\(f_{a}\) 是 \(\mathbf{Q}\) 到 \([0, a]\) 中的一个连续单射映射。试借助习题 9 推出：存在 \(\mathbf{Q}\) 到其自身的连续双射映射，其逆映射在任一点都不连续。

¶ 14) 设 X 是一个不只由一点组成的连通豪斯多夫拓扑空间，用 \(\Delta\) 表示 \(X \times X\) 的对角线。

\begin{enumerate}
\def\labelenumi{\alph{enumi})}
\item
试证：不存在把 \(\mathbb{C}\Delta\) 分成两个开集 \(\mathbf{D}_1, \mathbf{D}_2\) 的分划，使得 \(\mathbf{D}_1^{-1} = \mathbf{D}_1\) 且 \(\mathbf{D}_2^{-1} = \mathbf{D}_2\) 。{[}首先注意 \(\overline{\mathbf{D}_1(x)}\) 与 \(\overline{\mathbf{D}_2(x)}\) 对每个 \(x \in \mathbf{X}\) 都是连通的（利用第一章 § 11 习题 4 a)）；由此推出，若 \(y \in \mathbf{D}_1(x)\) ，则有 \(\overline{\mathbf{D}_2(x)} \times \{y\} \subset \mathbf{D}_1\) ，并证明这就蕴涵：对每个 \(x \in \mathbf{X}\) ，两集 \(\mathbf{D}_1(x)\) ， \(\mathbf{D}_2(x)\) 中必有一个为空；从而两集 \(\mathbf{D}_1, \mathbf{D}_2\) 中必有一个为空，这与假设矛盾。{]}由此推出：或者 \(\mathbb{C}\Delta\) 是连通的，或者 \(\mathbb{C}\Delta\) 恰有两个分支 A, B，使得 \(\mathbf{B} = \overline{\mathbf{A}}\) 。
\item
称 X 上的一个线性序与 X 的拓扑 \(\mathfrak{C}\) 相容，如果 \(\mathfrak{C}_0(\mathbf{X})\) 比 \(\mathfrak{C}\) 粗。试由 \(a\) ) 推出：X 上存在这样的线性序当且仅当 \(\mathbb{C}\Delta\) 不连通，并且这时恰有两个这样的线性序。{[}用 \(a\) ) 的记号，证明序关系必须是 \((x, y) \in \mathbf{A}\) 或 \((x, y) \in \mathbf{B} = \overline{\mathbf{A}}\) 。{]}试证：对这些序而言，各区间在拓扑 \(\mathfrak{C}\) 下都是连通的，并且它们是 X 关于 \(\mathfrak{C}\) 的仅有的连通子集。
\item
试证：若 X 在拓扑 \(\mathfrak{C}\) 下是局部连通或局部紧的，且 X 上存在一个与 \(\mathfrak{C}\) 相容的线性序，则必有 \(\mathfrak{C} = \mathfrak{C}_0(X)\) 。{[}若 X 局部连通，利用 b)。若 X 局部紧，通过考察邻域的边界，证明点 \(x \in X\) 的每个紧邻域也是 \(x\) 关于 \(\mathfrak{C}_0(X)\) 的邻域{]}。
\end{enumerate}

* d) 设 X 是 \(\mathbf{R}^2\) 的子空间，由 (o, o) 与所有点对 \((x, y)\) （满足 \(x > 0\) 且 \(y = \sin (\mathrm{i} / x)\) ）组成。试证：X 上存在一个线性序，它与 \(\mathbf{R}^2\) 的拓扑在 X 上诱导的拓扑相容，但该拓扑严格地细于 \(\mathcal{C}_0(X)\) 。 *

试给出一个连通豪斯多夫空间 X 的例子，它的任何一点都不具有连通邻域的基本系，而其上存在一个与 X 的拓扑相容的线性序{[}参看第一章，§ 11，习题 2 b){]}。

¶ 15) a) 设 X 是一个连通豪斯多夫空间，使得对每个 \(x \in X\) ，\(\complement\{x\}\) 恰有两个分支。试证：若 K 是这两个分支中的任何一个，则 \(K = K \cup \{x\}\) {[}第一章，§ 11，习题 4 a){]}。设 \(x, y\) 是 X 的两个不同的点，设 A 与 B 是 \(\complement\{x\}\) 的分支，并设 \(A', B'\) 是 \(\complement\{y\}\) 的分支。试证：两集 A, B 之一包含于 \(A'\) 或 \(B'\) 之中。

\begin{enumerate}
\def\labelenumi{\alph{enumi})}
\setcounter{enumi}{1}
\tightlist
\item
设 \(x_0\) 是 \(X\) 的一个点，并设 \(A(x_0), B(x_0)\) 是 \(\complement\{x_0\}\) 的分支。对每个 \(x \neq x_0\) ，\(\complement\{x\}\) 的两个分支中恰有一个或者包含于 \(A(x_0)\) 中，或者包含 \(A(x_0)\) ；把这个分支记作 \(A(x)\) 。试证：关系 \(A(x) \subset A(y)\) 是一个与 \(X\) 的拓扑相容的线性序关系（习题 14）。c) 把 b) 的结论推广到下述情形：\(\complement\{x\}\) 对每个 \(x \in X\) （除一个点 \(a \in X\) 外，对该点 \(\complement\{a\}\) 是连通的）都有两个分支。
\end{enumerate}

* \(d\) ) 设 \(\mathbf{X}\) 是 \(\mathbf{R}^2\) 的子空间，由点 \(a = (0, 1)\) ， \(b = (0, -1)\) 与所有点对 \((x, y)\) （满足 \(x \neq 0\) 且 \(y = \sin(1/x)\) ）组成。试证：\(\mathbf{X}\) 在拓扑 \(\mathfrak{G}\) （在 \(\mathbf{X}\) 上由 \(\mathbf{R}^2\) 的拓扑诱导）下是连通的，\(\mathbb{G}\{a\}\) 与 \(\mathbb{G}\{b\}\) 都是连通的，并且对每个点 \(x\) （ \(\mathbf{X}\) 中异于 \(a\) 与 \(b\) 的），\(\mathbb{G}\{x\}\) 恰有两个分支；但 \(\mathbf{X}\) 上不存在与拓扑 \(\mathfrak{G}_{*}\) 相容的线性序。

¶ 16) 设 X 是一个在两个不同点 \(a\) 与 \(b\) 之间完全不可约的连通豪斯多夫空间，也就是说，X 的连通子集中除 X 自身外没有一个同时含有 \(a\) 与 \(b\) 。试证：\(\mathbf{X}' = \mathbb{C}\{a, b\}\) 是连通的，并且对每个 \(x \in \mathbf{X}'\) ，\(x\) 在 X 中的余集恰有两个分支 \(\mathbf{A}(x)\) 与 \(\mathbf{B}(x)\) ，使得 \(a \in \mathbf{A}(x)\) ， \(b \in \mathbf{B}(x)\) ， \(a \notin \mathbf{B}(x)\) ， \(b \notin \mathbf{A}(x)\) 。{[}试利用第一章 § 11 习题 4 a) 证明 \(\mathbb{C}\{x\}\) 不能分成三个非空开集。{]}由此推出：X 上存在一个与 X 的拓扑相容的线性序（利用习题 15）。

¶ 17) 设 X 是一个连通豪斯多夫空间，使得每当 X 被三个异于 X 的非空连通集覆盖时，其中必有两个集的并异于 X。

\begin{enumerate}
\def\labelenumi{\alph{enumi})}
\item
试证：对每个点 \(x \in \mathbf{X}\) ，\(\mathbb{G}\{x\}\) 至多有两个分支（与习题 16 相同的方法）。
\item
试证：X 中至多有两个点，其余集是连通的。此外，若有两个不同的点 a, b 具有此性质，试证：对每个异于 a 或 b 的 x，a 与 b 属于 \(C\{x\}\) 的不同分支。
\item
试由 a) 与 b) 推出：X 上存在一个与 X 的拓扑相容的线性序（利用习题 15）。
\end{enumerate}

¶ 18) a) 设 X 是一个紧、连通、局部连通的空间，它在它的两个点 \(a\) , \(b\) 之间不可约（第二章，§ 4，习题 19）。设 \(x\) 是 X 中异于 \(a\) 或 \(b\) 的点。试证：\(\mathbb{C}\{x\}\) 恰有两个分支，且 \(a\) 与 \(b\) 属于 \(\mathbb{C}\{x\}\) 的不同分支。{[}按习题 16 的方式论证，首先证明 \(\mathbb{C}\{x\}\) 的分支不能多于两个。其次，若 V 是 \(x\) 的任一既不含 \(a\) 也不含 \(b\) 的连通闭邻域，令 \(\mathrm{A}_{\mathrm{V}}\) , \(\mathrm{B}_{\mathrm{V}}\) 分别是 \(a\) , \(b\) 在 \(\mathbb{C}\mathrm{V}\) 中的分支；利用第二章 § 4 习题 17 证明 \(\mathrm{A}_{\mathrm{V}} \neq \mathrm{B}_{\mathrm{V}}\) ；最后考察诸集 \(\mathrm{A}_{\mathrm{V}}\) （相应地 \(\mathrm{B}_{\mathrm{V}}\) ）当 V 取遍 \(x\) 在 X 中的所有连通闭邻域时的并。{]}由此推出：X 上存在一个线性序，使得 \(\mathfrak{C}_{0}(\mathbf{X})\) 就是 X 上给定的拓扑{[}习题 15 与 14 c){]}。

\begin{enumerate}
\def\labelenumi{\alph{enumi})}
\setcounter{enumi}{1}
\tightlist
\item
设 X 是一个紧连通空间，它在它的两个点 a, b 之间不可约。试证：若 X 的素成分（第二章，§ 4，习题 20）多于一个，则 X 的素成分所成的空间 \(X'\) （同处）在 a 与 b 的素成分之间不可约，并且 \(X'\) 上存在一个线性序，使得 \(X'\) 的拓扑等同于 \(\mathcal{G}_{0}(\mathbf{X}')\) 。 * 试给出一个 \(X'\) 异于 X 的例子{[}参看习题 15 d{]}。 *
\end{enumerate}

¶ 19) a) 设 X 是一个连通豪斯多夫空间，\(a, b\) 是 X 的两个不同的点，使得 X 的连通闭子集中除 X 自身外没有一个同时含有 \(a\) 与 \(b\) 。若 \(x\) 是 X 中异于 \(a\) 或 \(b\) 的点，则 \(\mathbb{C}\{x\}\) 不连通当且仅当：对每个 \(y \in \mathbb{C}\{x\}\) ，在 \(\mathbb{C}\{y\}\) 中存在一个在 X 中闭且含有 \(x\) 与两点 \(a, b\) 之一的连通子集。{[}为证条件必要，利用第一章 § 11 习题 4 a)；为证条件充分，令 A（相应地 B）是所有这样的 \(y \in X\) 所成之集：在 \(\mathbb{C}\{y\}\) 中存在一个在 X 中闭且含有 \(a\) 与 \(x\) （相应地 \(b\) 与 \(x\) ）的连通子集；证明 A 与 B 都非空、都是开集且互不相交，并且 \(\mathbb{C}\{x\} = A \cup B\) 。

\begin{enumerate}
\def\labelenumi{\alph{enumi})}
\setcounter{enumi}{1}
\tightlist
\item
设 X 是一个紧连通空间，它在它的两个点 \(a, b\) 之间不可约（第二章，§ 4，习题 19）。试证：若对 X 的每对不同的点 \(x, y\) ，X 中存在一个含有 \(y\) 与 \(a, b\) 之一但不含有 \(x\) 的紧连通集，则 X 上存在一个线性序，使得 \(\mathfrak{C}_{0}(\mathbf{X})\) 就是 X 上给定的拓扑（习题 15 与 16）。
\end{enumerate}

\begin{enumerate}
\def\labelenumi{\arabic{enumi})}
\setcounter{enumi}{19}
\item
在第一章 § 9 习题 23 的构造中，取 \(\mathbf{X}_0\) 为 R 的区间 {[}0, 1{]} 。试证：X 到其自身的同胚与 \(\mathbf{X}_0\) 到其自身的同胚相同，尽管 X 的拓扑严格地细于 \(\mathbf{X}_0\) 的拓扑{[}利用第一章 § 8 习题 20 b){]}。
\item
试证：若 \(r > 2\) ，则不存在 \(r\) 重传递的、\(\mathbf{R}\) 到其自身的同胚群（考察使 \(\mathbf{R}\) 的两点固定不动的子群）。
\end{enumerate}

\BourbakiExerciseGroup

\begin{enumerate}
\def\labelenumi{\arabic{enumi})}
\tightlist
\item
设 \(x\) 是一个实数 \(\geqslant 0\) ，并设 \(p\) 与 \(q\) 是两个整数 \(> 0\) 。试证下列关系式
\end{enumerate}

\[
(x ^ {1 / p}) ^ {1 / q} = x ^ {1 / p q}, \quad (x ^ {p}) ^ {1 / q} = (x ^ {1 / q}) ^ {p}, \quad x ^ {1 / p} x ^ {1 / q} = (x ^ {p + q}) ^ {1 / p q}.
\]

\begin{enumerate}
\def\labelenumi{\arabic{enumi})}
\setcounter{enumi}{1}
\tightlist
\item
  \begin{enumerate}
  \def\labelenumii{\alph{enumii})}
  \tightlist
  \item
在多项式环 \(\mathbf{A} = \mathbb{R}[\mathbf{X}]\) 上，考虑下述全序环结构：其中 \(\geqslant 0\) 的元素是 \(0\) 以及所有多项式 \(\neq 0\) （其首项系数 \(>0\) ）。试证：拓扑 \(\mathcal{G}_{0}(\mathbf{A})\) 与 \(\mathbf{A}\) 的环结构不相容。
  \end{enumerate}
\end{enumerate}

\begin{enumerate}
\def\labelenumi{\alph{enumi})}
\setcounter{enumi}{1}
\item
若 K 是一个有序域，则拓扑 \(\mathcal{G}_{0}(\mathbf{K})\) 与 K 的环结构相容，并且对这个环拓扑而言的有界集（第三章，§ 6，习题 12）就是对于序结构的有界集。拓扑 \(\mathcal{G}_{0}(\mathbf{K})\) 是局部逆有界的（第三章，§ 6，习题 22），特别地它与 K 的域结构相容；此外，K 关于这个拓扑的完备化 \(\hat{\mathbf{K}}\) （它是一个域，见同处）典范地具有有序域的结构{[}§ 1，习题 3 b){]}。
\item
在域 \(\mathbf{K} = \mathbf{Q}(\sqrt{2})\) （把它看作 R 的有序子域）上，考虑拓扑 C，它使得 \((x, y) \to x + y\sqrt{2}\) 是 \(Q^{2}\) 到 K 上的一个同胚，其中 \(Q^{2}\) 带有积拓扑。这个拓扑 C 细于 \(\mathcal{G}_{0}(\mathbf{K})\) 且与 K 的域结构相容；但这个拓扑域的完备化不是域。
\end{enumerate}

\(d\) ) 域 \(\mathbf{K} = \mathbf{Q}(\sqrt[3]{2})\) 作为 \(\mathbf{Q}\) 上的向量空间，是向量子空间 \(\mathbf{G}'\) （由 \(\mathbf{I}\) 与 \(\sqrt[3]{2}\) 生成）与向量子空间 \(\mathbf{G}''\) （由 \(\sqrt[3]{4}\) 生成）的直和。考虑 \(\mathbf{G}'\) 与 \(\mathbf{G}''\) 上由 \(\mathbf{R}\) 的拓扑诱导的拓扑；令 \(\mathfrak{C}\) 为 \(\mathbf{K}\) 上的、即 \(\mathbf{G}'\) 与 \(\mathbf{G}''\) 的拓扑之积的那个拓扑。试证：\(\mathfrak{C}\) 与 \(\mathbf{K}\) 的加法群结构相容，且细于 \(\mathfrak{C}_0(\mathbf{K})\) （ \(\mathbf{K}\) 被看作 \(\mathbf{R}\) 的有序子域），但不与 \(\mathbf{K}\) 的环结构相容。

¶ 3) a) 若 f 是域 R 到 R 的一个子域上的同构，试证 f 必是 R 到其自身的恒等映射。{[}注意所有有理数在 f 下不变，并且

\[
f (x ^ {2}) = (f (x)) ^ {2} \geqslant 0,
\]

从而 \(f\) 是递增的{]}。

\begin{enumerate}
\def\labelenumi{\alph{enumi})}
\setcounter{enumi}{1}
\tightlist
\item
设 \(K_{0}\) 是域 \(\mathbf{R}(\mathbf{X})\) （R 上一个未定元的有理分式域），其序这样规定：取多项式 \textgreater0 者为那些首项系数 \textgreater0 的多项式。设 K 为 \(K_{0}\) 的极大有序代数扩张。试证：K 的自同构中存在一个保序的 f，使得 \(f(\xi)=\xi\) 对一切 \(\xi\in\mathbb{R}\) 成立，且 \(f(\mathbf{X})=\mathbf{X}^{2}(*)\) 。
\end{enumerate}

\BourbakiExerciseGroup

\begin{enumerate}
\def\labelenumi{\arabic{enumi})}
\item
试证：\(\mathbf{R}\) 的每个有界半开区间同胚于 \(\mathbf{R}\) 的每个无界闭区间。若 \(a < b\) ，则三个区间 \(]a, b[, [a, b[, [a, b]]\) 中任何两个都互不同胚。
\item
试证：映射 \(x \to x / (1 + |x|)\) 是 \(\mathbf{R}\) 到开区间 \(] - 1, 1[\) 上的一个同胚。
\item
设 I 表示区间 \([0, 1]\) （ \(\mathbf{R}\) 中的），并设 \(f\) 是 I 到其自身的一个同胚，满足 \(f(0) = 0\) 与 \(f(1) = 1\) 。试证：存在一个连续映射 \(g\) ，它把 \(\mathbf{I} \times \mathbf{I}\) 映入 I，使得：(i) 对每个 \(x \in \mathbf{I}\) ， \(g(x, 0) = x\) 且 \(g(x, 1) = f(x)\) ；(ii) 对每个 \(y \in \mathbf{I}\) ，偏映射 \(x \to g(x, y)\) 是 I 到其自身的同胚，满足 \(g(0, y) = 0\) 与 \(g(1, y) = 1\) 。
\item
试证：\(\mathbf{R}\) 上由 \(\overline{\mathbf{R}}\) 的唯一一致结构诱导的一致结构严格地粗于 \(\mathbf{R}\) 的加法一致结构。
\item
当点 \((x, y)\) 趋于 \((+\infty, -\infty)\) （同时保持在 \(\mathbf{R} \times \mathbf{R}\) 中）时，试证函数 \(x + y\) 的聚点所成的集是 \(\overline{\mathbf{R}}\) 。
\end{enumerate}

同样地，当 \((x,y)\) 趋于 \((0,+\infty)\) （同时保持在 \(R\times R\) 中）时，xy 的聚点所成的集是 \(\overline{R}\) 。

\begin{enumerate}
\def\labelenumi{\arabic{enumi})}
\setcounter{enumi}{5}
\item
试证：每个有理函数 \(\mathbf{P}(x)/\mathbf{Q}(x)\) （系数在 R 中），只要它在所有点 x （ R 的、使 \(\mathbf{Q}(x)\neq0\) 的）处有定义，就可连续延拓到点 \(+\infty\) 与 \(-\infty\) （取值于 \(\overline{R}\) ）。
\item
设 X 是一个全序集。
\end{enumerate}

\begin{enumerate}
\def\labelenumi{\alph{enumi})}
\item
设 \(X_{1}\) 是 X 的完备化（集合论，第三章，§ 1，习题 15）。\(X_{1}\) 是一个全序集，其元素是 X 的这样的子集 A：(i) 若 \(x \in A\) 且 \(y \leqslant x\) ，则 \(y \in A\) ；(ii) 若 A 在 X 中有上确界，则该确界属于 A{[}这等价于说 A 在 \(\mathcal{G}_{0}(X)\) 中是闭的{]}；\(X_{1}\) 的元素按包含关系排序。\(X_{1}\) 关于拓扑 \(\mathcal{G}_{0}(X_{1})\) 是紧的{[}§ 2，习题 6 a){]}。映射 \(x \to ]\leftarrow,x]\) （ X 到 \(X_{1}\) 中的）是严格保序的，并且若把 X 等同于它在这个映射下的像，则 X 在 \(X_{1}\) 中稠密，且 \(\mathcal{G}_{0}(X)\) 由 \(\mathcal{G}_{0}(X_{1})\) 诱导。\(X_{1}\) 连通当且仅当 X 无间隙（§ 2，习题 7）。
\item
对每个势 c，试给出一个全序集 X 的例子，它在拓扑 \(\mathfrak{G}_{0}(\mathbf{X})\) 下连通，并且使得 X 的任何一点的任何基本邻域系都具有势 \(\geqslant c\) {[}利用 a) 与 §1 习题 4 b) 的方法{]}。
\item
用 \(a\) 的记号，令 \(\mathbf{X}_{2}\) 为字典积 \(\mathbf{X}_{1} \times \left\{-\mathbf{i}, 0, \mathbf{i}\right\}\) 的这样的子集：它是由下列诸点组成的集的余集：(i) 点 \((x, 0)\) ，其中 \(x \notin \mathbf{X}\) ；(ii) 点 \((x, 1)\) ，其中 \(x \in \mathbf{X}\) 且所有 \(y > x\) （ \(\mathbf{X}\) 中的）所成之集有最小元素；(iii) 点 \((x, -1)\) ，其中 \(x \in \mathbf{X}\) 且所有 \(y < x\) （ \(\mathbf{X}\) 中的）所成之集有最大元素。试证：\(\mathbf{X}_{2}\) 赋予拓扑 \(\mathfrak{C}_{0}(\mathbf{X}_{2})\) 后是紧且全不连通的，\(\mathbf{X}\) 可借助严格保序映射 \(x \to (x, 0)\) 等同于 \(\mathbf{X}_{2}\) 的一个稠密子集，并且 \(\mathfrak{C}_{0}(\mathbf{X}_{2})\) 在 \(\mathbf{X}\) 上诱导的拓扑是离散的。
\item
设 \(X'\) 是一个包含 X 的全序集，它在 X 上诱导给定的序结构，在拓扑 \(\mathcal{C}_{0}(\mathbf{X}')\) 下是紧的，并且 X 关于这个拓扑在 \(X'\) 中稠密。试证：存在一个连续保序的满射映射 \(f: X' \to X_{1}\) 与一个连续保序的满射映射 \(g: X_{2} \to X'\) ，两者都在 X 上诱导恒等映射。
\end{enumerate}

\BourbakiExerciseGroup

\begin{enumerate}
\def\labelenumi{\arabic{enumi})}
\item
设 \(X_{1}, X_{2}\) 是两个有向集，f 是定义在 \(X_{1} \times X_{2}\) 上的实值函数，使得：对每个 \(x_{1} \in X_{1}\) ，映射 \(x_{2} \to f(x_{1}, x_{2})\) 在 \(X_{2}\) 上递增；对每个 \(x_{2} \in X_{2}\) ，映射 \(x_{1} \rightarrow f(x_{1}, x_{2})\) 在 \(X_{1}\) 上递增。试证：f 在 \(\overline{R}\) 中关于 \(X_{1}\) 与 \(X_{2}\) 的截面滤子之积有极限，并且该极限是 f 的上确界。
\item
  \begin{enumerate}
  \def\labelenumii{\alph{enumii})}
  \tightlist
  \item
设 \(f\) 是定义在集合 \(\mathbf{X}\) 上的实值函数，\(\mathbf{A}\) 是 \(\mathbf{X}\) 的一个非空子集，\(\varphi\) 是定义在 \(f(\mathbf{A})\) 上的递增实值函数。若 \(\varphi\) 在点 \(a = \sup_{x \in \mathbf{A}} f(x)\) 处连续，则 \(\sup_{x \in \mathbf{A}} \varphi(f(x)) = \varphi(\sup_{x \in \mathbf{A}} f(x))\) 。
  \end{enumerate}
\end{enumerate}

\begin{enumerate}
\def\labelenumi{\alph{enumi})}
\setcounter{enumi}{1}
\tightlist
\item
设 f 是定义在由滤子 F 滤子化的集合 X 上的实值函数，\(\varphi\) 是定义在 f 关于 F 的聚点所成之集的一个开邻域 V 上的实值函数。试证：若 \(\varphi\) 在 V 上递增且连续，则
\end{enumerate}

\[
\lim \sup _ {\mathfrak {F}} (\varphi \circ f) = \varphi (\lim \sup _ {\mathfrak {F}} f).
\]

\begin{enumerate}
\def\labelenumi{\arabic{enumi})}
\setcounter{enumi}{2}
\item
设 \(f\) 是定义在无穷集 \(\mathbf{X}\) 上的实值函数，\(\mathfrak{G}\) 是 \(\mathbf{X}\) 的有限子集的余集所成的滤子。试证：\(\limsup_{\mathfrak{G}} f\) 是所有这样的实数 \(x\) （集合 \(f([x, +\infty])\) 为无穷集）所成之集的上确界。
\item
设 \(f, g\) 是定义在集合 \(X\) 上的两个实值函数，该集合由滤子 \(\mathfrak{F}\) 滤子化。试证
\end{enumerate}

\[
\lim \sup _ {\mathfrak {F}} (\sup (f, g)) = \sup (\lim \sup _ {\mathfrak {F}} f, \lim \sup _ {\mathfrak {F}} g).
\]

试给出一个由滤子 F 滤子化的集合 X，以及定义在 X 上的实值函数的一个无穷族 \((f_{i})\) ，使得

\[
\sup _ {t} (\lim \sup _ {\mathfrak {F}} f _ {t}) <   \lim \sup _ {\mathfrak {F}} (\sup _ {t} f _ {t}).
\]

\begin{enumerate}
\def\labelenumi{\arabic{enumi})}
\setcounter{enumi}{4}
\tightlist
\item
设 \(f, g\) 是定义在滤子化集合上的两个实值函数。试证：若 \(\lim g\) 存在且 \(\geqslant 0\) ，则
\end{enumerate}

\[
\lim \sup f g = (\lim \sup f) (\lim g)
\]

只要两边都有定义。

\begin{enumerate}
\def\labelenumi{\arabic{enumi})}
\setcounter{enumi}{5}
\item
设 \(X_{1}\) （相应地 \(X_{2}\) ）是由滤子 \(F_{1}\) （相应地 \(F_{2}\) ）滤子化的集合，f 是定义在 \(X_{1} \times X_{2}\) 上的实值函数。试用例子证明：三个数 \(\limsup_{\mathfrak{F}_{i} \times \mathfrak{F}_{2}} f(x_{1}, x_{2})\) ， \(\limsup_{\mathfrak{F}_{i}} (\limsup_{\mathfrak{F}_{2}} f(x_{1}, x_{2}))\) ， \(\limsup_{\mathfrak{F}_{i}} (\limsup_{\mathfrak{F}_{i}} f(x_{1}, x_{2}))\) 一般是不同的。
\item
设 \(f\) 是定义在闭子集 \(A\) （ \(\overline{\mathbf{R}}\) 的）上的实值函数。用 \(\limsup_{x \to a, x \geqslant a} f(x)\) {[}相应地 \(\limsup_{x \to a, x \leqslant a} f(x)\) {]}表示 \(f(x)\) 当 \(x\) 趋于点 \(a \in A\) 且保持在 \(A \cap [a, +\infty]\) 中时的上极限。
\end{enumerate}


（相应地 A ∩ {[}---∞, a{]}）。试证：点 \(a \in A\) （满足 \(\limsup_{x \to a, x \geqslant a} f(x) \neq \limsup_{x \to a, x \leqslant a} f(x)\) ）所成之集是可数的。{[}证明：对每对有理数 p, q （满足 p \textless{} q ），考虑集 \(G_{p,q}\) ，即点 \(a \in A\) 中使

\[
\lim _ {x \rightarrow a, x \geqslant a} \sup f (x) \leqslant p <   q \leqslant \lim _ {x \rightarrow a, x \leqslant a} \sup f (x)
\]

成立的那些点所成之集；它是可数的（利用 § 2 习题 1）{]}。

\begin{enumerate}
\def\labelenumi{\arabic{enumi})}
\setcounter{enumi}{7}
\item
试由习题 7 推出：实值函数 \(f\) ，它定义在闭子集 \(\mathbf{A}\) （ \(\overline{\mathbf{R}}\) 的）上，且在 \(\mathbf{A}\) 上单调，则除 \(\mathbf{A}\) 的一个可数子集的点外，它在 \(\mathbf{A}\) 上连续。
\item
设 X 是具有可数基的拓扑空间，f 是定义在 X 上的实值函数。试证：点 \(x \in X\) （使 \(\lim_{x \to a, x \neq a} f(x)\) 存在且异于 \(f(a)\) ）所成之集是可数的（与习题 7 相同的方法）。
\item
设 X 是具有可数基 \((\mathbf{U}_{n})\) 的拓扑空间，f 是定义在 X 上的实值函数；称 f 在点 \(a \in X\) 处达到严格相对极大值，如果存在 a 的一个邻域 V，使得 \(f(x) < f(a)\) 对 \(x \in V\) 中每个异于 x = a 的点成立。试证：X 中使 f 达到严格相对极大值的点所成的集 M 是可数的。（考察这样的 \(U_{n}\) 所成之集：f 在 \(U_{n}\) 的某点处达到等于其在 \(U_{n}\) 中上确界的严格相对极大值，并证明存在这个集到 M 上的映射）。
\item
设 \((u_n)\) 是实数 \(>0\) 的序列，满足 \(\lim_{n\to \infty}u_n = 0\) 。试证：存在无穷多个指标 \(n\) ，使得 \(u_{n}\geqslant u_{m}\) 对一切 \(m\geqslant n\) 成立。
\item
设 \((u_n)\) 是实数 \(>0\) 的序列，满足
\end{enumerate}

\[
\lim _ {n \to \infty} \inf u _ {n} = 0.
\]

试证：存在无穷多个指标 \(n\) ，使得 \(u_{n} \leqslant u_{m}\) 对一切 \(m \leqslant n\) 成立。

\begin{enumerate}
\def\labelenumi{\arabic{enumi})}
\setcounter{enumi}{12}
\item
设 \((u_{n})\) 是有限实数的序列，\((\varepsilon_{n})\) 是数 \(\geqslant 0\) 的序列，满足 \(\lim_{n\to \infty}\varepsilon_n = 0\) ，且 \(u_{n + 1}\geqslant u_n - \varepsilon_n\) 对一切整数 \(n\geqslant 0\) 成立。令 \(a = \liminf_{n\to \infty}u_n\) ， \(b = \limsup_{n\to \infty}u_n\) ；试证：序列 \((u_{n})\) 的聚点所成之集是区间 \([a,b]\) 。
\item
设 \((r_n)\) 是有限实数 \(>0\) 的递增序列，满足 \(\lim r_n = +\infty\) 。对每个有限实数 \(r > 0\) ，用 \(\mathbf{N}(r)\) 表示最大指标 \(n\) ，它满足 \(r_n \leqslant r\) 。试证
\end{enumerate}

\[
\limsup _ {r \to \infty} \frac {\mathrm{N} (r)}{r} = \limsup _ {n \to \infty} \frac {n}{r _ {n}}, \quad \liminf _ {r \to \infty} \frac {\mathrm{N} (r)}{r} = \liminf _ {n \to \infty} \frac {n}{r _ {n}}.
\]

¶ 15) 设 \((x_n)\) 是有限实数的序列，\((p_n)\) 是有限实数 \(\geqslant 0\) 的序列，满足 \(\lim_{n\to \infty}\left(\sum_{i=0}^{n}p_i\right) = +\infty\) 。令 \(y_n = \left(\sum_{i=0}^{n}p_ix_i\right) / \left(\sum_{i=0}^{n}p_i\right)\) ，这里 \(n\) 取使 \(\sum_{i=1}^{\infty}p_i \neq 0\) 的那些值。试证 \(\liminf_{n\to \infty}x_n \leqslant \liminf_{n\to \infty}y_n \leqslant \limsup_{n\to \infty}y_n \leqslant \limsup_{n\to \infty}x_n\) 。

设 H 是序列 \((x_{n})\) 在 \(\overline{R}\) 中的聚点所成之集的任一非空子集。试证：序列 \((p_{n})\) （有限实数 \(\geqslant0\) 的）可以这样确定，使得 \(\lim_{n\to\infty}\left(\sum_{i=0}^{n}p_{i}\right)=+\infty\) ，并且使相应的序列 \((y_{n})\) 的聚点所成之集包含 H。{[}化归到 H 可数的情形，然后用归纳法定义 \((p_{n})\) ，取其各项为 o 或 1{]}。

由此推出：序列 \((x_{n})\) 在 \(\overline{\mathbf{R}}\) 中收敛，当且仅当序列 \((y_{n})\) 在 \(\overline{\mathbf{R}}\) 中收敛——对每个序列 \((p_n)\) （其项为数 \(\geqslant 0\) 且满足 \(\lim_{n\to \infty}\left(\sum_{i = 0}^{n}p_i\right) = +\infty .\) ）皆然。

\begin{enumerate}
\def\labelenumi{\arabic{enumi})}
\setcounter{enumi}{15}
\tightlist
\item
设 \(x_0, y_0\) 是满足 \(0 < y_0 \leqslant x_0\) 的两个实数。对 \(n > 0\) ，用下列关系式递归地定义两个序列 \((x_n), (y_n)\)
\end{enumerate}

\[
x _ {n + 1} = \left(x _ {n} + y _ {n}\right) / 2, \quad y _ {n + 1} = \sqrt {x _ {n} y _ {n}}.
\]

试证：两个序列 \((x_{n})\) ， \((y_{n})\) 趋于同一极限 \(\alpha\) （即 \(x_{0}, y_{0}\) 的``算术-几何平均''）；此外，存在数 \(a > 0\) 与 \(\gamma\) ，使得 \(0 < \gamma < 1\) 且 \(x_{n} - y_{n} \leqslant a\gamma^{2n}\) 对每个 \(n\) 成立{[}注意 \(x_{n+1} - y_{n+1} = (x_{n} - y_{n})^{2}/4(x_{n+1} + y_{n+1})\) {]}。

\begin{enumerate}
\def\labelenumi{\arabic{enumi})}
\setcounter{enumi}{16}
\tightlist
\item
设 \(g\) 是 \(]0, i]\) 到 \([-1, 1]\) 中的一个映射，满足
\end{enumerate}

\[
\lim _ {x \rightarrow 0, x > 0} g (x) = 0.
\]

试证：存在一个连续递增映射 \(g_{2}\) 与一个连续递减映射 \(g_{1}\) ，它们把 {[}0, 1{]} 映入 {[}-1, 1{]} ，且满足 \(g_{1}(0) = g_{2}(0) = 0\) 。

并且 \(g_{1}(x) \leqslant g(x) \leqslant g_{2}(x)\) 对 \(0 < x \leqslant 1\) 成立。{[}对每个整数 \(n > 0\) ，考察下确界 \(x_{n}\) ——诸数 \(x\) （满足 \(g(x) \geqslant 1 / n\) ）的下确界{]}。

\begin{enumerate}
\def\labelenumi{\arabic{enumi})}
\setcounter{enumi}{17}
\tightlist
\item
把第 1 节到第 6 节的定义与结果推广到这样的函数：其值取在于拓扑 \(\mathfrak{G}_{0}(\mathbf{X})\) 下为紧的全序集 X 中（参看 § 2，习题 6）。
\end{enumerate}

\BourbakiExerciseGroup

\begin{enumerate}
\def\labelenumi{\arabic{enumi})}
\item
设 \(f\) 是开区间 \(\mathbf{I} \subset \mathbb{R}\) 到 \(\mathbb{R}\) 中的一个连续映射。试证：若 \(f(\mathbf{I})\) 是开集，且对每个 \(y \in \mathbb{R}\) ，集合 \(\overline{f}^1(y)\) 至多含两个不同的点，则 \(f\) 是单调的。
\item
设 B 是把 R 看作域 Q 上的向量空间时的一个基（``Hamel 基''）。B 是不可数的（§ 2，习题 2）。设 φ 是 B 的一个子集 C ≠ B 到集合 B 上的一个双射。如下定义 R 到 R 中的一个映射 f：若 \(x = \sum_{\xi \in B} \lambda(\xi) \xi\) ，其中 \(\lambda(\xi) \in \mathbf{Q}\) 且 \(\lambda(\xi) = 0\) 对除有限个元素 \(\xi \in B\) 外的一切元素成立，则 \(f(x) = \sum_{\xi \in C} \lambda(\xi) \varphi(\xi)\) 。试证：\(f(x + y) = f(x) + f(y)\) ，但对每个 \(z \in R\) ， \(f(z)\) 是 R 的一个稠密子集，由此推出 f 在 R 的任何区间上既不上有界也不下有界。
\item
  \begin{enumerate}
  \def\labelenumii{\alph{enumii})}
  \tightlist
  \item
对每个区间 \(\mathbf{I} \subset \mathbb{R}\) ，用 \(\mathbf{G}(\mathbf{I})\) 表示 \(\mathbf{I}\) 到其自身的所有同胚所成的群。试证：若 \(\mathbf{I}\) 与 \(\mathbf{J}\) 是 \(\mathbb{R}\) 的任意两个区间，每个都含多于一个点，则 \(\mathbf{G}(\mathbf{I})\) 与 \(\mathbf{G}(\mathbf{J})\) 同构。在 \(\mathbf{G}(\mathbf{I})\) 中，递增同胚所成的集 \(\mathbf{F}(\mathbf{I})\) 是指数为 2 的正规子群。
  \end{enumerate}
\end{enumerate}

\begin{enumerate}
\def\labelenumi{\alph{enumi})}
\setcounter{enumi}{1}
\item
设 \(\mathbf{G} = \mathbf{G}(\mathbf{R})\) ， \(\mathbf{F} = \mathbf{F}(\mathbf{R})\) 。对每个 \(f \in \mathbf{G}\) 与每个 \(x \in \mathbf{R}\) ，用 \(\sigma(x; f)\) 表示 \(\operatorname{sgn}(f(x) - x)\) 。设 \(\mathbf{T}^{+}\) （相应地 \(\mathbf{T}^{-}\) ）是所有这样的 \(f \in \mathbf{F}\) 所成之集：\(\sigma(x; f)\) 等于 \(\mathfrak{I}\) （相应地 \(-\mathfrak{I}\) ）且在 \(\mathbf{R}\) 上为常数。试证：\(\mathbf{T}^{+}\) （相应地 \(\mathbf{T}^{-}\) ）的每个元素在 \(\mathbf{F}\) 中共轭于平移 \(x \to x + \mathfrak{I}\) （相应地 \(x \to x - \mathfrak{I}\) ）。{[}若 \(f \in \mathbf{T}^{+}\) ，考察点列 \(f^{n}(0), n \in \mathbf{Z}\) {]}。
\item
设 \(T = T^{+} \cup T^{-}\) 。试证：每个不属于 T 的 \(f \in F\) 都是 \(T_{*}\) 中两个元素在 F 中的乘积{[}在 F 中写 \(f = f_{1}f_{2}\) ，其中 \(f_{1}(x) = \sup (x, f(x))\) ， \(f_{2}(x) = \inf (x, f(x))\) ；若 g 是平移 \(x \to x + 1\) ，则 \(f_{1}g\) 与 \(g^{-1}f_{2}\) 属于 T{]}。
\item
两个元素 \(f, g\) 在 \(F\) 中共轭，当且仅当存在 \(s \in F\) ，使得 \(\sigma(x; f) = \sigma(s(x); g)\) 对一切 \(x \in \mathbb{R}\) 成立。{[}考察连通分支 \(I_n\) ——即所有 \(x\) （满足 \(\sigma(x; f) \neq 0\) ）所成的开集的连通分支；注意 \(F(I_n)\) 与 \(F\) 同构（由 \(a\) ）并用 \(b\) {]}。特别地，若 \([a, b]\) 是 \(R\) 的一个区间，使得 \(\sigma(a; f) = \sigma(b; f) = 0\) 对某个 \(f \in F\) 成立，则如下元素 \(f'\) 属于 \(F\) ：它与 \(f\) 相重合（对 \(x \leqslant a\) 或 \(x \geqslant b\) ），且满足 \(f'(x) = x + \varepsilon(f(x) - x)\) （对 \(a < x < b\) ，其中 \(0 < \varepsilon < 1\) ）——这个元素与 \(f\) 在 \(F\) 中共轭。
\item
设 \(f \in F\) ，并设 \(a, b, c, d\) 是四个实数，满足 \(a < c < b < d\) ，且 \(f(x) - x = 0\) 对 \(x = a, x = b, x = c, x = d, \sigma(x; f) = 1\) 成立（于 \(a < x < c\) 与 \(b < x < d\) ）。设 \(a', b'\) 是两个数，满足 \(a \leqslant a' < c\) 与 \(b < b' \leqslant d\) ；令 \(J = [a, b], J' = [a', b']\) ，并用 \(f_1\) 表示 \(f\) 在 \(J\) 上的限制；于是 \(f_1 \in F(J)\) 。试证：存在一个递增同胚 \(s\) （ \(J\) 到 \(J'\) 上的），使得元素 \(f_2 = sf_1s^{-1}\) （ \(F(J')\) 的）满足：\(f_2(x) > f^{-1}(x)\) 每当 \(a' < x < b'\) {[}用 \(d\) {]}。
\end{enumerate}

\(f)\) 设 \(f \in \mathbf{F}\) ，并设 \([a, b]\) 是 \(\mathbf{R}\) 的一个区间，使得 \(f(x) - x = 0\) 对 \(x = a\) 与 \(x = b\) 成立，\(\sigma(x; f) = -1\) 对 \(x < a\) 成立，且 \(\sigma(x; f) = +1\) 对 \(x > b\) 成立。试证：存在 \(g \in \mathbf{F}\) ，它共轭于 \(f\) （在 \(\mathbf{F}\) 中），使得 \(g(x) > f(x)\) 对一切 \(x \in \mathbf{R}\) 成立（同样的方法）。

\(g\) ) 用 \(\mathbf{H}^{+}\) （相应地 \(\mathbf{H}^{-}\) ）表示 \(\mathbf{F}\) 的由所有这样的 \(f\in \mathbf{F}\) 组成的正规子群：\(f(x) = x\) 在 \(+\infty\) （相应地 \(-\infty\) ）的一个邻域内成立。试证：\(\mathbf{H}^{+},\mathbf{H}^{-}\) 与 \(\mathbf{H} = \mathbf{H}^{+}\cap \mathbf{H}^{-}\) 是 \(\mathbf{F}\) 的除 \(\mathbf{F}\) 与 \(\{e\}\) 外仅有的正规子群。{[}若 \(\mathbf{N}\) 是 \(\mathbf{F}\) 的一个正规子群，异于 \(\mathbf{F}\) ，且 \(f\in \mathbf{N}\) 不属于 \(\mathbf{H}^{+}\) ，试证存在一个元素 \(g\in \mathbf{N}\) ，使得所有 \(x\in \mathbf{R}\) （对其有 \(\sigma (x;g) = +1\) ）所成之集是一个区间 \([a, + \infty ]\) ，并且 \(g(x) = x\) 对 \(x\leqslant a\) 成立。为此，利用 \(e)\) 与 \(f)\) 的构造，以及 \(b)\) 与 \(c)\) 。要证明 \(\mathbf{N}\) 不包含除 \(\mathbf{H}\) 与 \(\{e\}\) 外的在 \(\mathbf{F}\) 中正规的子群，考察一个元素 \(f\in \mathbf{H}\) ，它异于 \(e\) ；设 \(a\) 与 \(b\) 是所有这样的 \(x\in \mathbf{R}\) （\(\sigma (x;f)\neq 0\) ）所成之集的下确界与上确界（由假设二者均有限）。然后注意 \(\mathrm{F}([a,b])\) 同构于 \(\mathbf{F}\) ，并用前面的结果{]}。

\begin{enumerate}
\def\labelenumi{\alph{enumi})}
\setcounter{enumi}{7}
\tightlist
\item
试证：群 H 是单群。（与证明 H 不含在 F 中正规的非平凡子群的方法相同）。
\end{enumerate}

\begin{enumerate}
\def\labelenumi{\arabic{enumi})}
\setcounter{enumi}{3}
\item
设 \(f\) 是拓扑空间 \(\mathbf{X}\) 上的一个下半连续函数，\(\varphi\) 是 \(f(\mathbf{E}) \subset \overline{\mathbf{R}}\) 上的一个递增下半连续函数。试证：\(\varphi \circ f\) 在 \(\mathbf{X}\) 上是下半连续的。
\item
设 \(f\) 是拓扑空间 \(\mathbf{X}\) 上的一个下半连续函数。试证：若 \(\mathbf{A}\) 是 \(\mathbf{X}\) 的任一非空子集，则 \(\sup f(\overline{\mathbf{A}}) = \sup f(\mathbf{A})\) 。
\item
设 \(f\) 是定义在拓扑空间 \(X\) 上的一个实值函数。数（有限或无穷）
\end{enumerate}

\[
\omega (a; f) = \lim _ {x \to a} \sup f (x) - \lim _ {x \to a} \inf f (x)
\]

称为 \(f\) 在点 \(a \in \mathbf{X}\) 处的振荡，只要右边有定义。

\begin{enumerate}
\def\labelenumi{\alph{enumi})}
\item
试证：\(x \to \omega(x; f)\) 在子空间 \(A\) （ \(X\) 的）上是上半连续的，该映射正是在其上有定义。
\item
若 \(f\) 在 \(\mathbf{X}\) 上有限，则 \(\mathbf{A} = \mathbf{X}\) ，且对每个 \(a \in \mathbf{X}\) 有
\end{enumerate}

\[
\omega (a; f) = \lim _ {(x, y) \rightarrow (a, a)} \sup (f (x) - f (y)).
\]

\begin{enumerate}
\def\labelenumi{\alph{enumi})}
\setcounter{enumi}{2}
\tightlist
\item
设 \(f\) 是 \(\mathbf{X}\) 上的一个下半连续有限实值函数。试证：若在一处 \(a \in \mathbf{X}, \omega(a; f)\) 取有限值，则
\end{enumerate}

\[
\liminf _ {x \to a} \omega (x; f) = 0.
\]

{[}设结论不真，试证将存在任意接近 a 的点 x，使 \(f(x)\) 可以任意大{]}。

\begin{enumerate}
\def\labelenumi{\alph{enumi})}
\setcounter{enumi}{3}
\tightlist
\item
对每个既约形式的有理数 \(r = p / q\) （其中 \(q > 0\) ），令 \(f(r) = q\) 。试证：\(f\) 在 \(\mathbf{Q}\) 上是下半连续的，但在每个点 \(r \in \mathbf{Q}\) 处有 \(\omega(r; f) = +\infty\) 。
\end{enumerate}

¶ 7) 设 X 是局部紧空间，f 是 X 上的一个下半连续函数。

\begin{enumerate}
\def\labelenumi{\alph{enumi})}
\tightlist
\item
若 \(\limsup_{x\to a}f(x)=+\infty\) 对每个 \(a\in\mathbf{X}\) 成立，试证：集 \(\overline{f}(+\infty)\) 在 \(\mathbf{X}\) 中稠密{[}对每个 \(a\in\mathbf{X}\) 与每个邻域 \(\mathbf{V}\) （ \(a\) 的），定义一个相对紧开集序列 \((\mathbf{U}_{n})\) ，使得 \(\overline{\mathbf{U}}_{n+1}\subset\mathbf{U}_{n}\) ，且 \(f(x)>n\) 对一切 \(x\in\mathbf{U}_{n}\) 成立{]}。
\end{enumerate}

* b) 设 \(n \to r_n\) 是 \(\mathbf{N}\) 到属于 {[}0, 1{]} 的有理数所成之集上的一个双射，并令 \(\varphi(x) = \mathrm{i}/\sqrt{|x|}\) {[}取 \(\varphi(0) = +\infty\) {]}；则函数 \(f(x) = \sum_{n=0}^{\infty} 2^{-n} \varphi(x - r_n)\) 在 {[}0, 1{]} 上是下半连续的，\(\overline{f}(+\infty)\) 在这个区间中稠密，其余集也如此{[}为证最后一点，注意以

\[
2 ^ {- n} \int_ {0} ^ {1} \varphi (x - r _ {n}) d x. *
\]

（积分学，第四章，§ 3，第 6 节，定理 5）。{]} *

\begin{enumerate}
\def\labelenumi{\alph{enumi})}
\setcounter{enumi}{2}
\item
若 \(f\) 有限，试证：点 \(x\) （使 \(\omega(x; f)\) 有限）所成之集是稠密的{[}用 a){]}。
\item
试证：X 中使 f 连续的点（无论 f 有限与否）所成之集是稠密的。{[}用 ( f/(1 + \textbar f\textbar) ) 代替 f（习题 4），化归到 f 有界的情形；注意 \(1/\omega(x; f)\) 是下半连续的，并用 a) 与习题 6c){]}。
\end{enumerate}

\begin{enumerate}
\def\labelenumi{\arabic{enumi})}
\setcounter{enumi}{7}
\item
设 X 是拓扑空间，A 是 \(X \times \overline{R}\) 的一个闭子集。试证：映射 \(x \to \inf(A(x))\) （ \(pr_{1}A\) 到 \(\overline{R}\) 中的）是下半连续的。反之，若 \(f: X \to \overline{R}\) 是下半连续函数，则 \(X \times \overline{R}\) 的子集 B（由对 \((x, y)\) （满足 \(y \geqslant f(x)\) ）组成）在 \(X \times \overline{R}\) 中是闭的。
\item
  \begin{enumerate}
  \def\labelenumii{\alph{enumii})}
  \tightlist
  \item
设 X 与 Y 是两个分离空间，\(\pi: X \to Y\) 是一个逆紧映射（第一章，§ 10，第 1 节）。设 \(g\) 是 X 上的一个下半连续实值函数。对每个 \(y \in Y\) ，令 \(f(y)\) 为 \(g\) 在集合 \(\overline{\pi}^{1}(y)\) 中的下确界{[}于是 \(f(y) = +\infty\) 若 \(\overline{\pi}^{1}(y) = \emptyset\) {]}。试证：\(f\) 在 Y 上是下半连续的。{[}利用如下事实：\(\pi(X)\) 是闭集，\(\overline{\pi}^{1}(y)\) 对一切 \(y \in \pi(X)\) 是紧的，以及 \(\overline{\pi}^{1}(y)\) 的每个邻域都包含一个关于等价关系 \(\pi(x) = \pi(x')\) 饱和的邻域。{]}
  \end{enumerate}
\end{enumerate}

\begin{enumerate}
\def\labelenumi{\alph{enumi})}
\setcounter{enumi}{1}
\tightlist
\item
设 \(g\) 是连续函数 \(|x_1 x_2 - 1|\) （定义在 \(\mathbf{R} \times ]0, +\infty[\) 上），对每个 \(x_1 \in \mathbb{R}\) 令 \(f(x_1) = \inf_{x_1 > 0} g(x_1, x_2)\) 。试证：\(f\) 不是下半连续的。
\end{enumerate}

\begin{enumerate}
\def\labelenumi{\arabic{enumi})}
\setcounter{enumi}{9}
\tightlist
\item
把定义 1、命题 1 与定理 3 推广到定义在拓扑空间 X 上、取值于任意全序集 Y 中的函数。若 f 与 g 是 X 到 Y 中的两个映射，都在一点 a 处下半连续，则 inf (f, g) 与 sup (f, g) 在 a 处下半连续。若 Y 是全序交换群，则 \(f + g\) 亦然。
\end{enumerate}

若 \(\mathbf{Y}\) 在拓扑 \(\mathfrak{G}_{0}(\mathbf{Y})\) （ \(\S 2\) ，习题 6）下是紧的，把定理 4 与命题 3、4 推广到取值于 \(\mathbf{Y}\) 中的函数。

¶ 11) 设 \(\overline{\mathbf{R}}\) 带有右拓扑（第一章，§ 1，习题 2）。一个映射 \(f\) （拓扑空间 \(\mathbf{X}\) 到 \(\overline{\mathbf{R}}\) 中的）在这个拓扑下连续，当且仅当 \(f\) 在 \(\mathbf{X}\) 的每一点处有相对极小。若 \(\mathbf{X} = \mathbf{R}\) （带通常拓扑），试证：集合 \(f(\mathbf{X})\) 于是是可数的。{[}设结论不真，且区间 \([\alpha, \beta] \cap f(\mathbf{X})\) 不可数；对每个 \(y \in [\alpha, \beta]\) ，令 \(\mathbf{U}(y) = \overline{f}([y, +\infty])\) ；\(\mathbf{U}(y)\) 是 \(\mathbf{R}\) 中的开子集。设 \((\mathbf{I}_n)\) 是 \(\mathbf{U}(\beta)\) 的连通分支的（有限或无穷）序列，对每个 \(n\) 令 \(x_n\) 为 \(\mathbf{I}_n\) 的一个点；对每个 \(y \in [\alpha, \beta]\) ，令 \(g_n(y)\) （相应地 \(h_n(y) \) ）为 \(\mathbf{U}(y)\) 的包含 \(x_n\) 的那个连通分支的左（相应地右）端点。试证：对至少一个 \(n\) 的值，函数 \(g_n, h_n\) 之一必在 \([\alpha, \beta]\) 中取不可数无穷多个不同的值；然后用 § 5 习题 8{]}。

\begin{enumerate}
\def\labelenumi{\arabic{enumi})}
\setcounter{enumi}{11}
\item
设 \(f\) 是紧区间 \([a, b]\) （ \(\mathbf{R}\) 的）上的一个连续、非常数的有限实值函数，满足 \(f(a) = f(b) = 0\) ；设 \(Z\) 是 \([a, b]\) 中由点 \(x\) （对其有 \(f(x) = 0\) ）组成的闭子集，\(c\) 是邻接于 \(Z\) 的诸区间在 \([a, b]\) 中长度的最大者。试证：对每个 \(t\) （满足 \(0 < t < c\) ），存在一点 \(x \in [a, b]\) ，使得 \(x + t \in [a, b]\) 且 \(f(x + t) = f(x)\) 。
\item
设 \(f\) 是紧区间 \([a, b]\) （ \(\mathbf{R}\) 的）上的一个连续有限实值函数，并设 \(f\) 不是严格单调的。试证：存在一点 \(x_0 \in ]a, b[\) ，使得对每个 \(\varepsilon > 0\) ，有两点 \(y, z\) 在 \([a, b]\) 中满足 \(x_0 - \varepsilon < y < x_0 < z < x_0 + \varepsilon\) 且 \(f(y) = f(z)\) 。
\item
设 \(f\) 是 \(\mathbf{R}\) 到其自身的一个连续映射。
\end{enumerate}

\begin{enumerate}
\def\labelenumi{\alph{enumi})}
\tightlist
\item
试证：若 \(f\) 在 \(\mathbf{R}\) 上一致连续，则存在两个实数 \(\alpha \geqslant 0\) 与 \(\beta \geqslant 0\) ，使得 \(|f(x)| \leqslant \alpha |x| + \beta\) 对一切 \(x \in \mathbf{R}\) 成立。b) 试证：若 \(f\) 在 \(\mathbf{R}\) 中单调且有界，则 \(f\) 在 \(\mathbf{R}\) 上一致连续。
\end{enumerate}

¶ 15) 设 X 是 Alexandroff 半直线{[}§ 2，习题 12 b){]}，并设 X′ 是它通过添加一个无穷远点 ω 所得的紧化。试证：对每个连续映射 f: X → X，或者我们有

\(\lim_{x \to \omega, x \in \mathbf{X}} f(x) = \omega,\) 或者存在 \(z \in \mathbf{X}\) 使得 \(f\) 为常数

在区间 \([z, \rightarrow[\) 中。{[}试证：若当 \(x\) 趋于 \(\omega, f\) 时有一个聚点 \(c \in \mathbf{X}\) ，则必有 \(\lim_{x \to \omega, x \in \mathbf{X}} = c\) ；证明 \(f\) 在 X′ 中不能有其他聚点，利用每个递增序列在 X 中收敛这一事实。然后用 § 2 习题 12 b) 以及在 X′ 中 ω 的任何可数个邻域之交仍是 ω 的邻域这一事实{]}。

\BourbakiExerciseGroup

\begin{enumerate}
\def\labelenumi{\arabic{enumi})}
\tightlist
\item
设 \((x_{t})\) 是一个实数族，其所有项都属于区间 \(A' = ] - \infty, +\infty]\) （相应地 \(A'' = [-\infty, +\infty[\) ）。族 \((x_{t})\) 在 \(\overline{\mathbf{R}}\) 中可和，当且仅当下列条件之一成立：
\end{enumerate}

\begin{enumerate}
\def\labelenumi{\alph{enumi})}
\item
数 \(x_{t}\) 中至少有一个等于 \(+\infty\) （相应地 \(-\infty\) ）。
\item
两个和 \(\sum_{i} x_{i}^{+}, \sum_{i} x_{i}^{-}\) 中至少有一个是有限的。
\end{enumerate}

在第一种情形有 \(\sum_{i} x_{i} = +\infty\) （相应地 \(-\infty\) ）；在第二种情形有 \(\sum_{i} x_{i} = \sum_{i} x_{i}^{+} - \sum_{i} x_{i}^{-}\) 。

\begin{enumerate}
\def\labelenumi{\arabic{enumi})}
\setcounter{enumi}{1}
\item
设 \((x_{i})\) 是一个实数族，其所有项都属于区间 \(\mathbf{A}'\) （相应地 \(\mathbf{A}''\) ），且满足习题 1 的条件 \(b\) )。试证：\((x_{i})\) 的每个子族在 \(\overline{\mathbf{R}}\) 中可和，并且 \(x'\) 的和是可结合的（参看第三章，§ 5，第 3 节，定理 2）。
\item
设 \((x_{n})_{n\geqslant 0}\) 是有限实数 \(\geqslant 0\) 的一个递减序列。这个序列在 \(\mathbf{R}\) 中可和，当且仅当序列 \((2^{n}x_{2^{n}})_{n\geqslant 0}\) 在 \(\mathbf{R}\) 中可和（``Cauchy 并项判别法''）。* 由此推出：通项为 \(1 / n^{\alpha}\) 的级数当 \(\alpha >1\) 时收敛，当 \(0 < \alpha \leqslant 1\) 时不收敛。*
\item
设 \((a_{n})\) 是有限实数 \(\geqslant 0\) 的任一序列。试证：序列 \(\left(\frac{a_n}{(1 + a_1)(1 + a_2)\ldots(1 + a_n)}\right)\) 在 \(\mathbf{R}\) 中可和（把这个序列的每一项表示为一个差）。
\item
设 \((d_n)\) 是有限实数 \(\geqslant 0\) 的一个序列，满足
\end{enumerate}

\[
\sum_ {n = 0} ^ {\infty} d _ {n} = + \infty .
\]

关于通项分别为下列各量的级数的收敛性，能说些什么？

\[
\frac {d _ {n}}{\mathrm{I} + d _ {n}}, \quad \frac {d _ {n}}{\mathrm{I} + n d _ {n}}, \quad \frac {d _ {n}}{\mathrm{I} + n ^ {2} d _ {n}}, \quad \frac {d _ {n}}{\mathrm{I} + d _ {n}}?
\]

\begin{enumerate}
\def\labelenumi{\arabic{enumi})}
\setcounter{enumi}{5}
\item
试证：\(s = \sum_{n=1}^{\infty} (1/n)\) 等于 \(+\infty\) ，办法是证明 \(s \geqslant s + \frac{1}{2}\) （求出偶指标各项之和与奇指标各项之和的、以 \(s\) 为函数的下界）。
\item
试证：对每个整数 \(p > \mathbf{r}\) ，有
\end{enumerate}

\[
\sum_ {n = 1} ^ {\infty} \frac {(- \mathrm{I}) ^ {n - 1}}{n ^ {p}} = (\mathrm{I} - 2 ^ {1 - p}) \sum_ {n = 1} ^ {\infty} \frac {\mathrm{I}}{n ^ {p}}.
\]

\begin{enumerate}
\def\labelenumi{\arabic{enumi})}
\setcounter{enumi}{7}
\item
试证：若 \(m\) 是一个整数 \(> 0\) ，则 \(\sum_{n \geqslant 1, n \neq m} \frac{1}{m^2 - n^2} = -\frac{3}{4m^2}\) （把有理分式 \(\frac{1}{m^2 - x^2}\) 表示为部分分式之和）。
\item
若通项为 \(x_{n}\) 的级数在 \(\mathbf{R}\) 中收敛，试证：\(\liminf_{n\to \infty}nx_n\leqslant 0\leqslant \limsup_{n\to \infty}nx_n\) 。
\end{enumerate}

¶ 10) 级数 \((u_n)\) 在 \(\mathbb{R}\) 中收敛，当且仅当：对每个递增序列 \((p_n)\) （数 \(>0\) ，满足 \(\lim_{n\to \infty}p_n = +\infty\) ），有 \(\lim_{n\to \infty}\left(\left(\sum_{k=0}^{n}p_ku_k\right) / p_n\right) = 0\) （用 § 5 习题 15）。

\begin{enumerate}
\def\labelenumi{\arabic{enumi})}
\setcounter{enumi}{10}
\tightlist
\item
考察一个序列 \((a_{n})\) ，它的每一项都是有限实数的一个有限序列之和
\end{enumerate}

\[
a _ {n} = b _ {n, 1} + b _ {n, 2} + \dots + b _ {n, k _ {n}}.
\]

对每对正整数 \((n, p)\) （满足 \(p \leqslant k_n\) ），若 \(m = \sum_{i=0}^{\infty} k_i + p\) ，令 \(c_m = b_{n,p}\) 。试证：若通项为 \(a_n\) 的级数在 \(\mathbf{R}\) 中收敛，且

\[
d _ {n} = \left| b _ {n, 1} \right| + \left| b _ {n, 2} \right| + \dots + \left| b _ {n, k _ {n}} \right|
\]

当 n 趋于无穷时趋于 o，则通项为 \(c_{m}\) 的级数收敛，且与通项为 \(a_{n}\) 的级数有相同的和。

¶ 12) 设 \((u_n)\) 是有限实数的一个序列。对集合 N 的每个置换 \(\sigma\) ，令

\[
r (n) = | \sigma (n) - n |. \sup _ {m \geqslant n} | u _ {m} |.
\]

\begin{enumerate}
\def\labelenumi{\alph{enumi})}
\item
若通项为 \(u_n\) 的级数收敛，且 \(\lim_{n \to \infty} r(n) = 0\) ，试证：通项为 \(v_n = u_{\sigma(n)}\) 的级数收敛到同一个和。{[}考察差 \(\sum_{k=0}^{n} v_k - \sum_{k=0}^{n} u_k\) （对大的 \(n\) 值），且若 \(h\) 是最小的整数 \(\leqslant n\) （它满足 \(\sigma(h) > n\) ），注意在这个差中，两个和的每一个里至多有 \(\sigma(h) - h\) 项不被消去{]}。
\item
设通项为 \(u_n\) 的级数收敛。试给出一个置换 \(\sigma\) 的例子，使得 \(\lim_{n \to \infty} r(n) = 0\) ，但不满足第三章，§ 5，习题 6 a) 的判别法。{[}定义一族两两不相交的区间 \(I_k = [n_k - 2k, n_k + 2k]\) （在 \(N\) 中），并取 \(\sigma\) 使得 \(\sigma(n) = n\) 每当 \(n\) 不属于任何 \(I_k\) 时成立，且 \(\sigma(n_k - 2j) = n_k + 2j\) ， \(\sigma(n_k + 2j) = n_k - 2j\) 对一切 \(k\) 与 \(0 \leqslant j \leqslant k\) 成立。{]}
\end{enumerate}

\begin{enumerate}
\def\labelenumi{\arabic{enumi})}
\setcounter{enumi}{12}
\item
设 \((u_n)\) 是实数 \(\geqslant 0\) 的一个序列，满足 \(\lim_{n\to \infty}u_n = 0\) 且 \(\sum_{n = 0}^{\infty}u_n = +\infty\) 。设 \((p_n)\) 是数 \(>0\) 的一个严格递增序列，满足 \(\lim_{n\to \infty}p_n = +\infty\) 。试证：存在一个置换 \(\sigma\) （ \(\mathbf{N}\) 的），使得对每个 \(n\in \mathbf{N}\) ， \(\sum_{k = 0}^{n}u_{\sigma (n)}\leqslant p_n\)
\item
设 \((u_{n})\) 是有限实数的一个序列，使得通项为 \(u_{n}\) 的级数收敛但不绝对收敛；令 \(s = \sum_{n=0}^{\infty} u_{n}\) 。试证：对每个数 \(s' \geqslant s\) ，存在一个置换 \(\sigma\) （ \(\mathbf{N}\) 的），使得 \(\sigma(n) = n\) 对那些指标 \(n\) （满足 \(u_{n} \geqslant 0\) ）成立，并且 \(\sum_{n=0}^{\infty} u_{\sigma(n)} = s'\) 。{[}对 \(m\) 用归纳法证明：存在一个置换 \(\sigma_m\) （ \(\mathbf{N}\) 的），使得 \(\sigma_m(k) = k\) 对一切 \(k\) （满足 \(u_k \geqslant 0\) ）成立，并且若用 \(u_k^{(m)}\) 表示 \(u_{\sigma_m(k)}\) ，则有一个整数 \(p_m\) 具有如下性质：\(\left| s' - \sum_{i=0}^{k} u_i^{(m)} \right| \leqslant 1/m\) 对一切 \(k \geqslant p_m\) 成立。此外，\(\sigma_{m+1}\) 满足：\(\sigma_{m+1}(k) = \sigma_m(k)\) 对一切 \(k\) （满足 \(\sigma_m(k) < p_m\) ）以及一切 \(k\) （满足 \(u_k \leqslant -1/m\) ）成立。{]}
\item
设 \((u_{n})\) 是有限实数的一个序列，满足 \(\lim_{n\to \infty}u_n = 0\) 且 \(\sum_{n}u_{n}^{+} = \sum_{n}u_{n}^{-} = +\infty\) 。若 \(a\) 与 \(b\) 是任意两个实数（有限或无穷），满足 \(a\leqslant b\) ，试证：存在一个置换 \(\sigma\) （ \(\mathbf{N}\) 的），使得当令 \(s_n = \sum_{k = 0}^n u_{\sigma (k)}\) 时，有 \(\liminf_{n\to \infty}s_n = a\) 与 \(\limsup_{n\to \infty}s_n = b\) 。试证：序列 \((s_n)\) 的聚点所成之集于是就是区间 \([a,b]\) 。
\item
设 \((u_n)\) 是有限实数的一个序列，使得通项为 \(u_n\) 的级数收敛但不绝对收敛。试证：存在一个置换 \(\sigma\) （ \(\mathbf{N}\) 的），它满足第三章，§ 5，习题 6 \(a\) ) 的条件，但使通项为 \(u_{\sigma - 1(n)}\) 的级数不收敛。{[}设 \(h, k, m\) 是具有下列性质的三个整数：若 \(p_1, \ldots, p_r\) 是整数 \(n\) （满足 \(h \leqslant n < h + m\) 且 \(u_n \geqslant 0\) ），则 \(u_{p_1} + \cdots + u_{p_r} > 2\) 且 \(h + m < k - r\) ；此外若 \(\mu, \nu\) 是使 \(h + m \leqslant \mu \leqslant \nu\) 的任意整数，则 \(\left|\sum_{n = \mu}^{\nu} u_n\right| \leqslant 1\) 。令 \(s = m - r\) ，并用 \(q_1, \ldots, q_s\) 表示这些整数 \(n\) （它们满足 \(h \leqslant n < h + m\) 且 \(u_n < 0\) ）。考察一个置换 \(\pi\) ——区间 \([h, k + s]\) （ \(\mathbf{N}\) 中的）的置换——它满足：\(\pi(p_i) = k - i + 1\) （对 \(i \leqslant i \leqslant r\) ）， \(\pi(q_j) = k + j\) （对 \(i \leqslant j \leqslant s\) ）， \(\pi(k + j) = h + s - j\) （对 \(i \leqslant j \leqslant s\) ）， \(\pi(k - i + 1) = h + s + i - 1\) （对 \(i \leqslant i \leqslant r\) ），最后 \(\pi(n) = n\) （对 \(h + m \leqslant n \leqslant k - r\) ）；试证我们有
\end{enumerate}

\[
\left| \sum_ {i = h} ^ {k} u _ {i} - \sum_ {i = h} ^ {k} u _ {\pi^ {- 1} (i)} \right| \geqslant I.
\]

\begin{enumerate}
\def\labelenumi{\arabic{enumi})}
\setcounter{enumi}{16}
\tightlist
\item
若 \(u_{mn} = \mathbf{i} / (m^2 - n^2)\) （对 \(m \neq n\) ）且 \(u_{mm} = 0\) ，试证
\end{enumerate}

\[
\sum_ {m = 1} ^ {\infty} \left(\sum_ {n = 1} ^ {\infty} u _ {m n}\right) = - \sum_ {n = 1} ^ {\infty} \left(\sum_ {m = 1} ^ {\infty} u _ {m n}\right) \neq 0.
\]

（用习题 8。）

\begin{enumerate}
\def\labelenumi{\arabic{enumi})}
\setcounter{enumi}{17}
\item
设 \((\mathbf{i} + u_{\mathfrak{i}})\) 是一个实数族，其所有项都属于区间 \([0, +\infty[\) （相应地 \(]0, +\infty]\) ）。族 \((\mathbf{i} + u_{\mathfrak{i}})\) 在 \(\mathbb{R}\) 中可乘，当且仅当下列条件之一成立：a) 数 \(\mathbf{i} + u_{\mathfrak{i}}\) 中至少有一个等于 o（相应地 \(+\infty\) ）。b) 族 \(u_{\mathfrak{i}}\) 在 \(\overline{\mathbb{R}}\) 中可和。
\item
设 \((\mathbf{I} + u_{t})_{t\in \mathbf{I}}\) 是 \(\mathbb{R}\) 中的一个可乘族。对每个非空子集 \(\mathbf{H}\) （ \(\mathbf{I}\) 的），用 \(v_{\mathbf{H}}\) 表示 \(\prod_{i\in \mathbf{H}}u_i\) 。试证：族 \((v_{\mathbf{H}})_{\mathbf{H}\in \mathfrak{F}(\mathbf{I})}\) 在 \(\mathbb{R}\) 中可和，并且
\end{enumerate}

\[
\mathrm{I} + \sum_ {\mathbf {H}} v _ {\mathbf {H}} = \prod_ {i \in \mathbf {I}} (\mathrm{I} + u _ {i}).
\]

逆命题是否成立？

由此推出：若 \(-\mathbf{I} < x < \mathbf{I}\) ，则 \(\prod_{n=0}^{\infty} (\mathbf{I} + x^{2^n}) = \mathbf{I} / (\mathbf{I} - x)\) 。

\begin{enumerate}
\def\labelenumi{\arabic{enumi})}
\setcounter{enumi}{19}
\item
设 \((u_n)\) 是有限实数 \(\geqslant 0\) 的一个序列，满足 \(u_0 > 0\) ，并令 \(s_n = \sum_{k=0}^{n} u_k\) （对每个 \(n \geqslant 0\) ）。则序列 \((u_n)\) 在 \(\mathbf{R}\) 中可和，当且仅当序列 \((u_n / s_n)\) 在 \(\mathbf{R}\) 中可和{[}对序列 \((1 - u_n / s_n)]\) 应用定理 4。若把序列 \((u_n / s_n)\) 换成序列 \((u_n / s_{n-1})\) ，结论同样成立。
\item
令 \(u_{n} = (-\mathrm{I})^{n} / \sqrt{n}\) （对 \(n \geqslant 2\) ）。以 \(\mathrm{I} + u_{n}\) 为通因子的乘积不收敛，但通项为 \(u_{n}\) 的级数收敛。
\item
对 \(n \geqslant 2\) ，令 \(u_{2n-1} = -\mathrm{I}/\sqrt{n}\) 且 \(u_{2n} = (\mathrm{I} + \sqrt{n})/n\) 。以 \(\mathrm{I} + u_n\) 为通因子的乘积收敛，但通项为 \(u_n\) 的级数不收敛。
\end{enumerate}

\BourbakiExerciseGroup

\begin{enumerate}
\def\labelenumi{\arabic{enumi})}
\tightlist
\item
若 \(x\) 与 \(y\) 是实数，我们有
\end{enumerate}

\[
[ x + y ] = [ x ] + [ y ] + \varepsilon \quad \text { where } \varepsilon = 0 \quad \text { or } \quad \varepsilon = 1,
\]

\[
[ x - y ] = [ x ] - [ y ] - \varepsilon \quad \text { where } \varepsilon = 0 \quad \text { or } \quad \varepsilon = 1,
\]

\[
[ x ] + [ x + y ] + [ y ] \leqslant [ 2 x ] + [ 2 y ],
\]

\[
[ x ] + \left[ x + \frac {1}{n} \right] + \left[ x + \frac {2}{n} \right] + \dots + \left[ x + \frac {n - 1}{n} \right] = [ n x ],
\]

\[
\left[ \frac {[ n x ]}{n} \right] = [ x ]
\]

对任何整数 \(n > 0\) 成立。

\begin{enumerate}
\def\labelenumi{\arabic{enumi})}
\setcounter{enumi}{1}
\item
设 \((\varepsilon_{n})\) 是趋于 0 的有限数 \(>0\) 的一个严格递减序列。对每个 \(x \in \mathbb{R}\) ，令 \(r_{n}(x)\) 为 \(\varepsilon_{n}\) 的那个倍数，它是 \(x\) 的精确到 \(\varepsilon_{n}\) 的不足近似。则序列 \((r_{n}(x))\) 对一切 \(x \in \mathbb{R}\) 递增，当且仅当 \(\varepsilon_{n}\) 是 \(\varepsilon_{n+1}\) 的整数倍（对每个 \(n\) ）。
\item
设 \(a\) 是一个整数 \(> 1\) 。实数 \(x\) 是有理数，当且仅当它以 \(a\) 为底的展开式，比如说 \(x = p_0 + \sum_{n=0}^{\infty} u_n a^{-n}\) ，是周期的，也就是说，当且仅当存在两个整数 \(n_0\) 与 \(r > 0\) ，使得 \(u_{n+r} = u_n\) 对一切 \(n \geqslant n_0\) 成立。
\item
设 \((u_{n})\) 是数 \(>0\) 的一个可和序列（在 \(\mathbf{R}\) 中），满足下列条件：\(u_{n+1} \leqslant u_{n}\) 且 \(u_{n} \leqslant \sum_{k=1}^{\infty} u_{n+k}\) 对一切 \(n \geqslant 0\) 成立。试证：对每个数 \(a\) （满足 \(0 < a \leqslant s = \sum_{n=0}^{\infty} u_{n}\) ），存在 \(\mathbf{N}\) 的一个子集 I，使得 \(a = \sum_{n \in \mathbf{I}} u_{n}\) 。若 \(u_{n+1} \leqslant u_{n} \leqslant 2u_{n+1}\) 对一切 \(n\) 成立，这些条件特别地被满足。考察二进展开的情形。
\item
对每个实数 \(x \in ]0, I]\) ，存在唯一的无穷递增序列 \((q_n)\) （整数 \(> 0\) 的），使得 \(x = \sum_{n=1}^{\infty} I / (q_1 q_2 \ldots q_n)\) 。数 \(x\) 是有理数，当且仅当 \(q_{n+1} = q_n\) 从 \(n\) 的某个值起成立。
\item
对每个实数 \(x > \mathbf{I}\) ，存在唯一的无穷序列 \((t_n)\) （整数 \(\geqslant \mathbf{I}\) 的），使得 \(t_{n+1} \geqslant t_n^2\) 对一切 \(n \geqslant \mathbf{I}\) 成立，且使得 \(x = \prod_{n=1}^{\infty} \left( \mathbf{I} + \frac{\mathbf{I}}{t_n} \right)\) （用 §7 习题 19）。数 \(x\) 是有理数，当且仅当 \(t_{n+1} = t_n^2\) 从 \(n\) 的某个值起成立。{[}为证条件是必要的，试证：若令 \(x_n = \prod_{k=n+1}^{\infty} \left( \mathbf{I} + \frac{\mathbf{I}}{t_k} \right)\) ，则分式 \(\frac{x_n}{x_n - \mathbf{I}}\) （取既约形式）的分母构成一个递减序列{]}。
\end{enumerate}

¶ * 7) 设 \((\varepsilon_n)_{n \geqslant 0}\) 是一个每项为 1 或 --- 1 的数的序列。试证：实数

\[
x _ {n} = \varepsilon_ {0} \sqrt {2 + \varepsilon_ {1} \sqrt {2 + \varepsilon_ {2} \sqrt {2 + \cdots + \varepsilon_ {n} \sqrt {2}}}}
\]

有定义且等于

\[
2 \sin \left(\frac {\pi}{4} \sum_ {k = 0} ^ {n} \frac {\varepsilon_ {0} \varepsilon_ {1} \dots \varepsilon_ {k}}{2 ^ {k}}\right).
\]

由此推出：\(\lim_{n\to \infty}x_n\) 存在，且对每个实数 \(x\) （满足 \(-2\leqslant x\leqslant 2\) ），存在一个序列 \((\varepsilon_n)\) ，使得 \(x\) 等于相应序列 \((x_{n})\) 的极限。对哪些 \(x\) 的值序列 \((\varepsilon_n)\) 是唯一的？对哪些 \(x\) 的值它是周期的？（见习题 3）。*

\begin{enumerate}
\def\labelenumi{\arabic{enumi})}
\setcounter{enumi}{7}
\item
试证：不存在非常数映射 \(\psi\) （区间 \(\mathbf{I} \subset \mathbf{R}\) 到自然数序列全体所成之集 \(\mathbf{N}^{\mathbb{N}}\) 中的），它对每个 \(x \in \mathbf{I}\) 具有下列性质：对每个整数 \(n \geqslant 0\) ，存在一个邻域 \(\mathbf{V}\) （ \(x\) 在 \(\mathbf{I}\) 中的），使得对每个 \(y \in \mathbf{V}\) ，前 \(n\) 项（序列 \(\psi(y)\) 的）与前 \(n\) 项（序列 \(\psi(x)\) 的）相同（注意：这将蕴涵 \(\psi\) 的连续性，若每个因子 \(\mathbf{N}\) 带有离散拓扑而 \(\mathbf{N}^{\mathbb{N}}\) 带有积拓扑）。
\item
试证：Cantor 三分集 K（§ 2，第 5 节）是所有这样的实数 \(x \in [0, 1]\) 所成之集，它们的三进展开（相应地，若 \(x\) 是 K 的一个邻接区间的端点，则为非正常三进展开）
\end{enumerate}

\[
x = \sum_ {n = 1} ^ {\infty} u _ {n} 3 ^ {- n}
\]

满足 \(u_{n} \neq \mathrm{I}\) （对每个 \(n\) ，因而 \(u_{n} = 0\) 或 \(u_{n} = 2\) ）。若 \(v_{n} = \frac{\mathrm{I}}{2} u_{n}\) 对每个 \(n\) 成立，且 \(f(x) = \sum_{n=1}^{\infty} v_{n} 2^{-n}\) ，试证：\(f\) 是 \(\mathbf{K}\) 到 \([\mathrm{o}, \mathrm{I}]\) 上的一个满射连续映射；由此推出 \(\mathbf{K}\) 具有连续统的势。

¶ 10) 设 \((d_n)\) 是一个基序列，令 \(a_n = d_n / d_{n-1}\) （ \(n \geqslant 1\) ），设 \((n_i)\) 是整数的一个严格递增序列，并设 \((b_i)\) 是整数的一个序列，满足 \(0 < b_i < a_{n_i} - 1\) （对每个 \(i\) ）。设 \(G\) 为所有这样的 \(x \in [0, 1]\) 所成之集：若 \(\sum_{n=1}^{\infty} u_n / d_n\) 是 \(x\) 的展开式（或非正常展开式，若它存在），则 \(u_{n_i} \neq b_i\) 对一切 \(i \in N\) 成立。试证：\(G\) 是一个全不连通完备集；并且若 \(l\) 是邻接于 \(G\) 且含于 \([0, 1]\) 的诸区间长度之和，则

\[
\mathrm{I} - l = \prod_ {i = 0} ^ {\infty} \left(\mathrm{I} - \frac {\mathrm{I}}{a _ {n _ {i}}}\right).
\]

¶ 11) 设 X 是一个豪斯多夫拓扑空间。假设对每个各项为 o 或 1 的有限序列 s，存在 X 的一个非空子集 \(\mathrm{A}(s)\) ，满足下列条件：(i) 若 s 是 n 项（每项为 o 或 i）的序列，而 \(s'\) , \(s''\) 是 \(n + 1\) 项（每项为 o 或 i）且前 n 项与 s 的相同的序列，则 \(\mathrm{A}(s) = \mathrm{A}(s') \cup \mathrm{A}(s'')\) ；又若 \(s_0\) 是空序列，则 \(\mathrm{A}(s_0) = \mathrm{X}\) 。

\begin{enumerate}
\def\labelenumi{(\roman{enumi})}
\setcounter{enumi}{1}
\tightlist
\item
对每个各项为 0 或 1 的无穷序列 \((u_n)_{n\geqslant 0}\) ，若用 \(s_n\) 表示有限序列 \((u_k)_{0\leqslant k\leqslant n}\) ，则由诸 \(\mathbf{A}(s_n)\) 形成的滤子基收敛于 \(\mathbf{X}\) 的一点。
\end{enumerate}

试证：在这些条件下，存在 Cantor 三分集 K 到 X 上的一个满射连续映射。

¶ 12) 试由习题 11 推出：

\begin{enumerate}
\def\labelenumi{\alph{enumi})}
\item
若 A 是 R 的一个紧子集，则存在 Cantor 三分集 K 到 A 上的一个连续映射{[}取 A(s) 为 A 与适当选取的区间的交{]}。
\item
若此外 A 还是完备且全不连通的，则它同胚于 K{[}同样的方法，只需这样安排：若 \(s_{1}\) 与 \(s_{2}\) 是两个有相同项数（各项为 o 或 1）的不同序列，则
\end{enumerate}

\[
\mathbf {A} (s _ {1}) \cap \mathbf {A} (s _ {2}) = \varnothing_ {1} ].
\]

¶ 13) 设 X 是一个可数的豪斯多夫空间。假设对每个各项为 o 或 1 的有限序列 s，存在 X 的一个非空子集 B(s)，满足下列条件：

\begin{enumerate}
\def\labelenumi{(\roman{enumi})}
\item
若 \(s\) 是 \(n\) 项（等于 \(\mathbf{o}\) 或 \(\mathbf{i}\) ）的有限序列，且 \(s', s''\) 是 \(n + \mathbf{i}\) 项（等于 \(\mathbf{o}\) 或 \(\mathbf{i}\) ）而前 \(n\) 项与 \(s\) 的相同的序列，则 \(\mathbf{B}(s') \cup \mathbf{B}(s'') = \mathbf{B}(s)\) 且 \(\mathbf{B}(s') \cap \mathbf{B}(s'') = \emptyset\) ；又若 \(s_0\) 是空序列，则 \(\mathbf{B}(s_0) = \mathbf{X}\) 。
\item
对每个 \(x \in \mathbf{X}\) ，若 \(s_n\) 是 \(n\) 项（等于 0 或 1）的（唯一的）序列，它满足 \(x \in \mathrm{B}(s_n)\) ，则由诸 \(\mathrm{B}(s_n)\) 形成的滤子基收敛于 \(x\) （在 \(\mathbf{X}\) 中）。
\end{enumerate}

试证：在这些条件下，E 同胚于有理直线 Q。{[}首先证明 X 同胚于 Cantor 三分集 K 的一个可数子集，它在 K 中稠密且不含 K 的任何邻接区间的端点。为此，注意假设使每个 \(x \in X\) 对应一个各项等于 0 或 1 的无穷序列 \((u_{n}(x))\) ；利用 X 可数这一事实，证明双射 \(s \to \mathbf{B}(s)\) 可以这样修改，使得对任何 \(x \in E\) ，诸 \(u_{n}(x)\) 都不构成最终常数的序列；最后用 §2 的习题 10{]}。

由此推出：\(\mathbf{R}\) 的每个没有孤立点的可数子空间同胚于 \(\mathbf{Q}\) 。

\begin{enumerate}
\def\labelenumi{\arabic{enumi})}
\setcounter{enumi}{13}
\tightlist
\item
试证：\(\mathbf{R}\) 的每个闭子集或者是可数的，或者具有连续统的势{[}用上面的习题 12 b) 以及第一章，§ 9 的习题 16{]}。
\end{enumerate}


\begin{enumerate}
\def\labelenumi{\arabic{enumi})}
\setcounter{enumi}{14}
\item
试证：\(\mathbf{R}\) 的开子集所成之集与 \(\mathbf{R}\) 的紧、全不连通、完备子集所成之集都具有连续统的势。
\item
  \begin{enumerate}
  \def\labelenumii{\alph{enumii})}
  \tightlist
  \item
设 A 是 R 的一个可数闭子集，f 是定义在 R 上的一个连续实值函数。试证：若 f 在邻接于 A 的每个区间中是常数，则 f 是常数（用 Bolzano 定理）。
  \end{enumerate}
\end{enumerate}

\begin{enumerate}
\def\labelenumi{\alph{enumi})}
\setcounter{enumi}{1}
\tightlist
\item
若 \(f\) 是习题 9 中定义的 Cantor 三分集 \(\mathbf{K}\) 到 \([0, 1]\) 上的那个连续映射，试证：\(f\) 可以延拓为 \(\mathbf{R}\) 上的一个连续函数，它在邻接于 \(\mathbf{K}\) 的每个区间上是常数。
\end{enumerate}

\begin{enumerate}
\def\labelenumi{\arabic{enumi})}
\setcounter{enumi}{16}
\tightlist
\item
设 X 是有一个可数稠密子集的拓扑空间。试证：X 上所有实值连续函数所成之集具有连续统的势。
\end{enumerate}
